\documentclass[10pt,a4paper]{amsart}
\usepackage[margin=19mm]{geometry}
\usepackage{amsmath,amssymb,mathtools}
\usepackage[scr=rsfs,cal=euler]{mathalfa}
\usepackage{xfrac}
\usepackage[unicode,hidelinks,bookmarks=false]{hyperref}
\newcommand{\ie}{i.e.\ }
\newcommand{\cf}{cf.\ }
\newcommand{\resp}{respectively\ }
\newcommand{\Subsection}{\subsection}
\providecommand{\nobreakdash}{\nobreak\hbox{-}}
\setlength{\emergencystretch}{2em}
\multlinegap0pt
\usepackage{amssymb}
\newtheorem{proposition}{Proposition}
\title{Boundedness of the Hardy-Littlewood maximal operator on generalized fofana spaces}
\author{Pokou Nagacy, Berenger Akon Kpata,, Nouffou Diarra}
\date{Source: arXiv:2605.21670v1 (CC BY 4.0)}
\begin{document}
\maketitle

\noindent\emph{Source and licence.} Boundedness of the Hardy-Littlewood maximal operator on generalized fofana spaces. Version arXiv:2605.21670v1 (CC BY 4.0), distributed by arXiv under the Creative Commons Attribution 4.0 International licence, http://creativecommons.org/licenses/by/4.0/, as recorded in the arXiv licence metadata for this exact version and shown on the abstract page https://arxiv.org/abs/2605.21670v1. Full licence notice: \texttt{licenses/arxiv-2605.21670-CC-BY-4.0.txt}.\par\smallskip

\noindent\textit{Editorial scope.} Counted unit: the article's proposition proposition (source0.tex lines 296--299). The article's complete original proof of it is delivered with it (source0.tex lines 300--346, one \texttt{proof} environment of 47 lines). No mathematics is written, repaired or re-derived: the statement and the proof are the article's own lines, in the article's own notation, and no retained line is substituted or re-flowed.  Retained uncounted as the source-local context the proof needs: lines 267--275.  Nothing was omitted from any retained span, so the package carries no omission record.  No retained line contains a citation, so the wrapper carries no bibliography and no bibliography entry.  No retained line contains a figure environment, an image-inclusion command or drawing code, so no image is delivered and the package declares no asset dependency.  Editorial formatting only, in addition: an AMS \texttt{amsart} preamble that restates the article's own declarations and macros used by the retained lines, namely the environments \texttt{proposition} on separate counters, together with the standard packages those lines need.  The retained environments print in this document as Proposition 1 in order of appearance, whereas the article prints the same items with the numbering of its own full document.  The complete native source of the article as distributed in the arXiv e-print for this exact version is bundled unchanged as \texttt{sources/201075/source0.tex}.
 \begin{proposition}\label{prop 1}
 Let  $ f \in L^{0}(\mathbb{R}^{d}). $    
 \begin{enumerate}
 \item If $1\leq q_{1}\leq q_{2} \leq p \leq\infty$ and $ \phi \in \mathcal{G}_{q_{2},p} $ then 
 $$\left\| f \right\|_{(L^{q_{1}},L^{p})^{\phi }(\mathbb{R}^{d})} \lesssim \left\|f\right\|_{(L^{q_{2}},L^{p})^{\phi }(\mathbb{R}^{d})},$$ and consequently, $(L^{q_{2}},L^{p})^{\phi}(\mathbb{R}^{d})\subset (L^{q_{1}},L^{p})^{\phi}(\mathbb{R}^{d}).$
 \item If $1\leq q\leq p_{1} \leq p_{2} \leq\infty$ and $ \phi \in \mathcal{G}_{q,p_{1}} $ then 
 $$\left\| f \right\|_{(L^{q},L^{p_{2}})^{\phi}(\mathbb{R}^{d})} \lesssim \left\|f\right\|_{(L^{q},L^{p_{1}})^{\phi}(\mathbb{R}^{d})}$$ and consequently, $(L^{q},L^{p_{1}})^{\phi}(\mathbb{R}^{d})\subset (L^{q},L^{p_{2}})^{\phi}(\mathbb{R}^{d}).$
 \end{enumerate}
 \end{proposition}

\begin{proposition}
 Let $1\leq q \leq  p \leq \infty$ and let $ \phi \in \mathcal{G}_{q,p}. $ Then 
 $$\frac{1}{\phi(r_{0})} \approx \left\| \chi_{B(0,r_{0})} \right\|_{(L^{q},L^{p})^{\phi }(\mathbb{R}^{d})}, $$ for every ball $ B(0,r_{0}). $
 \end{proposition}

\begin{proof}
Let $ r_{0} >0. $ It follows from Proposition \ref{prop 1} (2) that 
\begin{eqnarray*}
 \frac{1}{\phi(r_{0})} r_{0}^{-\frac{d}{q}} \left[\int_{\mathbb R^{d}} \chi_{B(0,r_{0})}(x)\chi_{B(0,r_{0})}(x)dx\right]^{\frac{1}{q}} &\leq &         \left\|\chi_{B(0,r_{0})} \right\|_{(L^{q},L^{+\infty})^{\phi}(\mathbb{R}^{d})}\\
 & \lesssim & \left\|\chi_{B(0,r_{0})}\right\|_{(L^{q},L^{p})^{\phi}(\mathbb{R}^{d})}.
\end{eqnarray*}
Thus, $ \frac{1}{\phi(r_{0})} \lesssim \left\| \chi_{B(0,r_{0})} \right\|_{(L^{q},L^{p})^{\phi }(\mathbb{R}^{d})}. $\\
We will prove the converse inequality in two steps. By definition of $\left\| \cdot \right\|_{(L^{q},L^{p})^{\phi }(\mathbb{R}^{d})},$ we have 
\begin{eqnarray*}
&&\left\| \chi_{B(0,r_{0})} \right\|_{(L^{q},L^{p})^{\phi }(\mathbb{R}^{d})} \\&=& \sup_{r>0} \frac{1}{\phi(r)} r^{d(-\frac{1}{q}-\frac{1}{p})} \left\|\left[\int_{\mathbb R^{d}} \chi_{B(0,r_{0})}(x)\chi_{B(\cdot,r)}(x)dx\right]^{\frac{1}{q}}\right\|_p. 
\end{eqnarray*}
\textit{First step:} For any real number $r$ such that $ r_{0} \leq r  $, we have
\begin{equation} \label{eqq 1}
r_{0}^{\frac{d}{q}}\phi(r_{0}) \leq C_{1} r^{\frac{d}{q}}\phi(r).
\end{equation}
It follows that
\begin{eqnarray*}
&&\sup_{r \ge r_{0}} \frac{1}{\phi(r)} r^{d(-\frac{1}{q}-\frac{1}{p})} \left\|\left[\int_{\mathbb R^{d}}  \chi_{B(0,r_{0})}(x)\chi_{B(\cdot,r)}(x)dx\right]^{\frac{1}{q}}\right\|_p \\&\leq& \sup_{r \ge r_{0}} \frac{1}{\phi(r)} r^{d(-\frac{1}{q}-\frac{1}{p})} \left\|\left[\int_{\mathbb R^{d}}  \chi_{B(0,r_{0})}(x)dx\right]^{\frac{1}{q}}\chi_{B(0,r+r_{0})}\right\|_p \\&\leq& 
 2^{d(\frac{1}{q}+\frac{1}{p})}
r^{\frac{d}{q}}_{0}\sup_{r \ge r_{0}} \frac{1}{\phi(r)} r^{d(-\frac{1}{q}-\frac{1}{p})} (r+r_{0})^{\frac{d}{p}}\\
&\leq& 2^{d(\frac{1}{q}+\frac{1}{p})}r^{\frac{d}{q}}_{0}\sup_{r \ge r_{0}} \frac{1}{\phi(r)} r^{d(-\frac{1}{q}-\frac{1}{p})} r^{\frac{d}{p}}(1+\frac{r_{0}}{r})^{\frac{d}{p}} \\
&\leq& 2^{d(\frac{1}{q}+\frac{2}{p})}r^{\frac{d}{q}}_{0}\sup_{r \ge r_{0}} \frac{1}{\phi(r)} r^{d(-\frac{1}{q}-\frac{1}{p})} r^{\frac{d}{p}}.
\end{eqnarray*} 
By inequality (\ref{eqq 1}), we get 
\begin{equation}\label{eq 2}
\sup_{r \ge r_{0}} \frac{1}{\phi(r)} r^{d(-\frac{1}{q}-\frac{1}{p})} \left\|\left[\int_{\mathbb R^{d}}  \chi_{B(0,r_{0})}(x)\chi_{B(\cdot,r)}(x)dx\right]^{\frac{1}{q}}\right\|_p \leq \frac{C^{'}}{\phi(r_{0})},
\end{equation}
where $C^{'}= 2^{d(\frac{1}{q}+\frac{2}{p})}C_{1}.  $\\
\textit{Second step:} For any real number $r$ such that $ r_{0} \geq r$, we have
\begin{equation}\label{eq 4}
r_{0}^{\frac{d}{p}}\phi(r_{0}) \leq C_{2} r^{\frac{d}{p}}\phi(r).
\end{equation}

\begin{eqnarray*}
&&\sup_{r \leq r_{0}} \frac{1}{\phi(r)} r^{d(-\frac{1}{q}-\frac{1}{p})} \left\|\left[\int_{\mathbb R^{d}}  \chi_{B(0,r_{0})}(x)\chi_{B(\cdot,r)}(x)dx\right]^{\frac{1}{q}}\right\|_p \\&\leq& \sup_{r \leq r_{0}} \frac{1}{\phi(r)} r^{d(-\frac{1}{q}-\frac{1}{p})} \left\|\left[\int_{\mathbb R^{d}}  \chi_{B(y,r)}(x)dx\right]^{\frac{1}{q}}\chi_{B(0,r+r_{0})}\right\|_p. \\
 &\leq& 2^{d(\frac{1}{q}+\frac{1}{p})}
\sup_{r \leq r_{0}} \frac{1}{\phi(r)} r^{d(-\frac{1}{q}-\frac{1}{p})}r^{\frac{d}{q}} (r+r_{0})^{\frac{d}{p}}\\
&\leq& 2^{d(\frac{1}{q}+\frac{1}{p})}r_{0}^{\frac{d}{p}}\sup_{r \leq r_{0}} \frac{1}{\phi(r)} r^{d(-\frac{1}{p})} (1+\frac{r}{r_{0}})^{\frac{d}{p}} \\
&\leq& 2^{d(\frac{1}{q}+\frac{2}{p})}r^{\frac{d}{p}}_{0}\sup_{r \leq r_{0}} \frac{1}{\phi(r)} r^{-\frac{d}{p}}.
\end{eqnarray*} 
By inequality (\ref{eq 4}), we get 
\begin{equation}\label{eq 3}
\sup_{r \ge r_{0}} \frac{1}{\phi(r)} r^{d(-\frac{1}{q}-\frac{1}{p})} \left\|\left[\int_{\mathbb R^{d}}  \chi_{B(0,r_{0})}(x)\chi_{B(\cdot,r)}(x)dx\right]^{\frac{1}{q}}\right\|_p \leq \frac{C^{''}}{\phi(r_{0})}
\end{equation}
where $C^{''}= 2^{d(\frac{1}{q}+\frac{2}{p})}C_{2}.  $\\
From (\ref{eq 2}) and (\ref{eq 3}), we deduce that $$\left\| \chi_{B(0,r_{0})} \right\|_{(L^{q},L^{p})^{\phi }(\mathbb{R}^{d})} \leq \frac{C}{\phi(r_{0})},$$ where $C=\max\{C^{'},  C^{''}  \}$.
\end{proof}

\end{document}
